% !TeX root = ../main.tex
\section{Experiments}
\label{sec:experiments}

\subsection{Protocol and Evidence Scope}
\label{sec:exp-protocol}
Captioning compares SFT checkpoint 1,000 with six audio-language
baselines~\cite{qwen3omni,mossaudio,mossmusic,musicflamingo,audioflamingo3,qwen25omni}.
All return 2,000 texts; training/checkpoint provenance needs a final
audit before claiming unseen-data generalization. BLEU-4 uses 0--100;
other agreement scores are multiplied by 100. BERTScore uses
RoBERTa-large, layer 17, without IDF or baseline rescaling, on the same
1,861 valid cases for every method; other caption metrics use all 2,000.
Style evaluation uses the 822-clip population only. Automatic caption
ratings use separate versioned protocols. Historical fixed-window
controllers differ from the shared-evidence bar protocol.

\subsection{Short-Clip Caption Agreement}
\label{sec:exp-reference}
SFT leads all six agreement measures (Table~\ref{tab:e2-full2000-v2}),
reaching 9.74 BLEU-4, 30.80 ROUGE-L, and 89.71 BERTScore F1 versus
the strongest baseline values of 0.85, 14.99, and 82.85. This consistency
supports task adaptation; shared references prevent interpreting these
metrics as independent music-quality evidence.

Motion-keyword precision/recall reach 76.23/85.90 versus Qwen3-Omni's
35.09/24.67, showing both broader reference-category coverage and a
lower proportion of non-reference categories. The lexicons omit
laterality, negation, action--body binding, and temporal correctness.
Temporal-connective presence scarcely separates Qwen2.5-Omni
(1,998/2,000) from SFT (2,000/2,000), motivating stricter evidence
requirements. Table~\ref{tab:e2-full2000-diagnostics} retains surface
diagnostics without treating duplication or length as dance quality.

\subsection{Automatic Caption Assessment}
\label{sec:exp-refinement}
Table~\ref{tab:refinement-d123} reports D1--D3 ratings of our captioning
pipeline. Ours obtains 4.97, 3.33, and 4.80 under r4, and 4.92, 4.50,
and 4.91 under v2.1. Means use valid ratings separately for each
dimension; D4--D6 remain unreported.

Both protocols score the same captions, so differences between rows
reflect scoring sensitivity. D1 is near its ceiling: 98.2\% of valid
r4 judgments assign 5. These results describe our system's ratings;
baseline comparisons and independent musical validation remain pending.
Appendix~\ref{sec:experimental-details} documents coverage, the inference
configuration, and the unverified mapping of source rating IDs to our
proposed anchors.

\subsection{Audio-Conditioned Style Prior}
\label{sec:exp-jev}
Table~\ref{tab:style-prior-822} compares three supervised MuQ-JEV
heads with seven zero-shot audio-language models on the same 822 clips.
Rankings use 22 labels, with 13 observed classes and positive training
examples in 15 classes. Focal reaches 44.77 Top-1, 81.39 Top-3,
32.17 Macro-F1, and 32.35 balanced accuracy, versus 22.38, 54.50,
11.21, and 14.15 for the strongest zero-shot comparator, Qwen2.5-Omni.
All three training objectives exceed that comparator on every measure.

Cross-entropy has the highest Top-3 (81.87); focal has the highest
point estimates on the other three metrics. Focal's paired Top-$k$
intervals against cross-entropy include zero. Weighted cross-entropy
trades Top-$k$ accuracy for slightly higher class-balanced scores than
ordinary cross-entropy. These results support supervised style
specialization against video-VLM pseudo-labels. Scarce classes and
incomplete training support limit coverage; human music--dance
suitability and downstream harness benefits require separate evidence.

\subsection{Controller Reliability and Composition}
\label{sec:exp-harness}
% !TeX root = ../main.tex
% Sources: formal_result_v1.json, judge_size_robustness.json,
% extension_set.json in experiments/structured_dev_20261114.
% The extension's attempted=9 and n_paired=8 supersede its stale n=6 field.
\begin{table}[t]
\centering
\caption{Preliminary fixed-window whole-song comparison. Completion reports W/O over all attempts; $n$ counts completed pairs. Entries show O$-$W mean total-score differences and 95\% song-bootstrap intervals. Judges see acoustic summaries and text, not audio; both belong to the planner's Qwen family.}
\label{tab:harness-preliminary}
\begingroup\footnotesize
\setlength{\tabcolsep}{2pt}
\renewcommand{\arraystretch}{1.12}
\begin{tabular*}{\columnwidth}{@{\extracolsep{\fill}}lccrr@{}}
\toprule
Cohort & Completion & $n$ & Qwen 8B & Qwen 30B \\
\midrule
Original & 6/11; 11/11 & 6 & +2.50 & +2.42 \\
 & & & $[0.08,3.92]$ & $[0.00,4.25]$ \\
Extension & 8/9; 9/9 & 8 & +6.00 & +4.44 \\
 & & & $[5.12,6.88]$ & $[3.38,5.56]$ \\
\bottomrule
\end{tabular*}
\endgroup
\end{table}

Only Workflow and Self-Refine complete all 24 songs in the initial
four-policy study (Table~\ref{tab:harness-controller-history}).
Self-Refine uses nearly three times Workflow's calls; tool-loop
policies cost more and complete fewer songs. Additional iteration alone
does not ensure delivery. The old quality check gives all 87 evaluable
plans its maximum, preventing quality discrimination.

The later W/O study uses Qwen3-VL-30B-A3B-Instruct at temperature zero,
approximately 5.06~s windows, and no JEV prior. O adds acoustic guidance
and bounded rewriting with unequal evidence access, a system-level
contrast. Table~\ref{tab:harness-preliminary} separates 11 development
songs from nine subsequent FMA Electronic songs. Completion counts
all attempts; scores use six/eight completed pairs. O completes all
extension songs, whereas W misses one.

Qwen judges favor O by +6.00/+4.44 out of 20 on the extension after
averaging orders. Both see text and acoustic summaries, not audio,
and share the planner's family. Cross-family text-only MOSS yields
$+0.50\;[-1.00,2.00]$, leaving independent musical confirmation
unresolved. These historical scores cannot validate the new rubric
or bar/style mechanisms.

\subsection{Exploratory Mechanism Analysis}
\label{sec:exp-mechanisms}
The repaired four-song comparison (Table~\ref{tab:harness-exploratory})
favors acoustically guided variants under its text-judge setup;
the quota-only interval crosses zero. Unequal evidence and development
reuse preclude isolated component attribution.

Early repair reduces mean normalized exact-text repetition from 73.14\%
to 36.82\% on five valid pairs of six registered songs, all decreasing.
This optimized diagnostic does not establish movement diversity.
Appendix~\ref{sec:remaining-tests} reports presentation sensitivity and
further controls.
